\documentclass[11pt]{article}

\usepackage[a4paper,margin=1in]{geometry}
\usepackage{amsmath,amssymb}
\usepackage{booktabs}
\usepackage{graphicx}
\usepackage{hyperref}
\usepackage{multirow}
\usepackage{enumitem}

\title{Compact Local Classification with Ranked Pattern-Structure Classifier Candidates}
\author{Alan Tomat \and TODO: Student Coauthors}
\date{\today}

\begin{document}
\maketitle

\begin{abstract}
Lazy classification with pattern structures avoids constructing the full
concept lattice, but it can still generate a large number of query-dependent
classifier candidates for each query object. Each candidate is obtained from a
query-source intersection and has an interpretable pattern description, but
aggregating a large unranked set can reduce both predictive quality and
explanation clarity. This paper studies a ranking-and-selection framework for
LazyFCA-style classification. For each query, candidate classifiers generated
for all classes are placed in a common pool, ranked by candidate-importance
metrics, and only the top $k$ candidates are retained for voting. We evaluate
whether compact ranked candidate sets preserve or improve predictive
performance while greatly reducing explanation size. The study compares
purity-based, support-based, locality-based, statistical, and hybrid ranking
metrics on tabular numerical datasets. Preliminary motivation suggests that
locality-aware purity metrics, such as query-weighted precision and
query-weighted log-odds ratio, are promising for compact local explanations.
\end{abstract}

\section{Introduction}

Formal Concept Analysis (FCA) provides a mathematical framework for deriving
interpretable concepts from binary object-attribute data. Pattern structures
extend FCA to complex descriptions such as intervals, sets, sequences, and
graphs. This makes it possible to use FCA-style concepts in machine learning
settings where objects have numerical, categorical, or structured attributes.

A central obstacle is scalability. Constructing a full concept lattice can be
expensive or infeasible for realistic datasets. LazyFCA addresses this by
classifying each query object locally. Instead of building a full lattice, it
intersects the query with training objects and obtains a set of local
candidate classifiers. Each candidate has an interpretable pattern description
and can be evaluated against same-class and opposite-class training objects.

However, a second problem appears: LazyFCA may generate many candidates per
query. Using all of them equally can be harmful. Some candidates are highly
specific, redundant, weakly supported, or misleading with respect to the query.
Thus, while each individual pattern description is interpretable, the final
explanation can become too large, and the voting procedure can be noisy.

This paper investigates whether candidate ranking can make LazyFCA both more
compact and more accurate. The guiding question is:

\begin{quote}
Can we rank local pattern-structure classifier candidates so that a small
top-$k$ subset
retains most of the predictive quality and yields a more interpretable local
explanation?
\end{quote}

The main contributions are intended to be:

\begin{itemize}[leftmargin=*]
    \item A global pooled top-$k$ selection framework for LazyFCA classifier
    candidates, where candidates from all classes compete in one ranking.
    \item A systematic empirical comparison of candidate-importance metrics
    for compact LazyFCA classification.
    \item An analysis of the tradeoff between predictive quality and number of
    retained candidates.
    \item A discussion of singleton source-query candidates, where $tp=1$ is
    generated by construction and should not be interpreted as independent
    multi-object support.
\end{itemize}

\section{Background}

\subsection{Formal Concept Analysis}

Formal Concept Analysis (FCA) represents data by a formal context
$(G,M,I)$, where $G$ is a set of objects, $M$ is a set of attributes, and
$I \subseteq G \times M$ records which objects have which attributes. FCA was
introduced as a framework for organizing knowledge through hierarchies of
concepts \cite{wille2009restructuring}; the standard mathematical treatment is
given by Ganter and Wille \cite{ganter1999formal}.

The basic construction is a pair of derivation operators. For a set of objects
$A \subseteq G$, the common attributes of those objects are
\[
    A' = \{m \in M \mid (g,m) \in I \text{ for all } g \in A\}.
\]
For a set of attributes $B \subseteq M$, the objects sharing all attributes in
$B$ are
\[
    B' = \{g \in G \mid (g,m) \in I \text{ for all } m \in B\}.
\]
A formal concept is a pair $(A,B)$ such that $A'=B$ and $B'=A$. The set
$A$ is the extent of the concept, and $B$ is its intent. Thus a concept is not
only a cluster of objects and not only a rule-like description: it is a maximal
set of objects together with the maximal set of attributes common to exactly
those objects.

Formal concepts are ordered by extent inclusion, equivalently by reverse intent
inclusion. This order forms the concept lattice. The lattice order has a
natural interpretation: concepts with larger extents and smaller intents are
more general, while concepts with smaller extents and richer intents are more
specific. This explicit object-description duality is the main reason FCA is
useful for interpretable data analysis: a prediction or candidate description can be
inspected through the objects it covers and the attributes that justify the
coverage \cite{carpineto2004concept,ignatov2014introduction}.

\subsection{Pattern Structures}

Classical FCA assumes that object descriptions are binary attributes. Pattern
structures extend the same concept-forming idea to complex descriptions
\cite{ganter2001pattern,kuznetsov2009pattern}. A pattern structure is usually
written as
\[
    (G,(D,\sqcap),\delta),
\]
where $G$ is a set of objects, $D$ is a description space, $\sqcap$ is a meet
or similarity operation on descriptions, and $\delta:G \to D$ maps each object
to its description. The description space is ordered by the meet operation:
we write $d_1 \sqsubseteq d_2$ when $d_1 \sqcap d_2 = d_1$. With this
convention, $d_1$ is at least as general as $d_2$.

The derivation operators are analogous to those in FCA. For
$A \subseteq G$, the common description of the objects in $A$ is
\[
    A^{\diamond} = \mathop{\sqcap}\limits_{g \in A} \delta(g).
\]
For a description $d \in D$, the objects described by $d$ are
\[
    d^{\diamond} = \{g \in G \mid d \sqsubseteq \delta(g)\}.
\]
A pattern concept is a pair $(A,d)$ such that $A^{\diamond}=d$ and
$d^{\diamond}=A$. The formal role of the intent is therefore played by a
structured description rather than by a set of binary attributes.

This generalization is important for machine-learning data because numerical,
set-valued, sequential, relational, or graph-like objects need not be reduced
to a large binary table before concept-based analysis is possible
\cite{kuznetsov2013fitting,kaytoue2015pattern,buzmakov2016mining}. Different
choices of $D$ and $\sqcap$ lead to different kinds of interpretable
descriptions. In this paper, the main description language is the interval
pattern structure for numerical data.

\subsection{Interval Pattern Structures}

Interval pattern structures describe numerical objects by axis-parallel
intervals. For a $d$-dimensional numerical object
$x=(x_1,\ldots,x_d)$, the object description can be viewed as the degenerate
interval vector
\[
    \delta(x) = ([x_1,x_1],\ldots,[x_d,x_d]).
\]
A general interval description is a vector
\[
    d = ([\ell_1,u_1],\ldots,[\ell_d,u_d]),
\]
where each component gives the admissible range of one feature. The natural
coverage relation is coordinate-wise: an object $x$ is covered by $d$ when
$\ell_j \leq x_j \leq u_j$ for every feature $j$. Thus an interval pattern is
readable as a hyper-rectangle in the original feature space.

The meet operation used here is the component-wise common generalization of
interval descriptions. For two one-dimensional intervals
$[\ell_1,u_1]$ and $[\ell_2,u_2]$, this operation returns
\[
    [\min(\ell_1,\ell_2), \max(u_1,u_2)].
\]
For several numerical features, the operation is applied independently to
each coordinate. Under this convention, wider intervals are more general
because they cover more objects, and narrower intervals are more specific.
This is the interval-pattern analogue of replacing several binary attributes
by their common description. Numerical pattern mining with FCA and pattern
structures has been studied as an alternative to losing information through
coarse binarization \cite{article}.

For the local classifier candidates studied in this paper, each candidate is
generated from a query object $x$ and a source training object $s$. Their
common interval description is

\[
    [\min(x_j, s_j), \max(x_j, s_j)]
\]

for each numerical feature $j$. Across all numerical features, this is an
axis-parallel hyper-rectangle containing both objects. This construction is
local to the query: a different query object produces a different family of
interval candidates even when the training set is fixed. In the main
numerical experiments, the retained explanation for a prediction is therefore
a small set of such query-dependent interval patterns. When mixed data are
processed by the implementation, categorical variables are one-hot encoded and
the categorical part of a candidate is represented by shared active Boolean
features; however, the paper-level comparison is intended to focus on
numerical datasets where the interval description is the primary explanatory
object.

\subsection{LazyFCA Classification}

Constructing a full concept or pattern-concept lattice may be unnecessary or
too expensive when the task is to classify one query object. LazyFCA-style
classification avoids full lattice construction by generating candidate
classifiers only around the query. Lazy classification with interval pattern
structures was
studied for credit scoring by Masyutin, Kashnitsky, and Kuznetsov
\cite{masyutin2015lazy}; related interval-pattern lazy classifiers and local
interval explanations are also considered in recent work on IPS-based
classification \cite{TOMAT2026109754}.

The terminology requires some care. In concept-based learning with pattern
structures, a positive hypothesis has a stricter formal meaning: it is a
pattern that is the common description of some positive examples and does not
cover negative examples \cite{kuznetsov2009pattern,kuznetsov2013fitting}.
Masyutin et al. distinguish $\alpha$-weak premises, also called classifiers,
from $\alpha$-weak hypotheses; the former may be more general than a common
description of positive examples, while the latter require equality with such a
common description \cite{masyutin2015lazy}. Therefore, in this paper a
query-source intersection is not called a hypothesis by default. It is a
candidate classifier, or premise candidate, until its coverage and ranking are
evaluated.

Let the training set be divided into classes. For a query object $q$ and a
class $c$, every training object $s$ with label $c$ is used as a source object.
The meet of the query description and the source description produces a local
candidate classifier
\[
    z(q,s,c).
\]
For interval data, this candidate is the smallest axis-parallel interval
pattern containing both $q$ and $s$. The candidate is then evaluated against
the training data: same-class objects are supporters, and objects from all
other classes are opposers. A candidate is useful when it covers enough
supporters while avoiding opposers, although different LazyFCA variants define
and aggregate usefulness in different ways.

For each generated candidate $z$ we define:

\begin{align*}
tp(z) &= \text{number of class-}c\text{ training objects covered by }z,\\
fp(z) &= \text{number of non-class-}c\text{ training objects covered by }z,\\
fn(z) &= \text{number of class-}c\text{ training objects not covered by }z,\\
tn(z) &= \text{number of non-class-}c\text{ training objects not covered by }z.
\end{align*}

In binary LazyFCA, this procedure is often described in terms of positive and
negative premises or classifiers. In the multiclass setting used in this paper, the same
logic is applied one class at a time: for each class $c$, the supporters are
the training objects of class $c$, and the opposers are all remaining training
objects. All generated class-specific candidates can then be aggregated to
predict the class of $q$.

A practical difficulty is that the number of local candidates can be large,
often on the order of the number of training objects. Moreover, each candidate
is built from the query and a same-class source object, so the source object is
covered by construction and $tp(z) \geq 1$. Consequently, a high-purity
candidate may sometimes be a narrow source-query description rather than a
broadly supported formal hypothesis or rule. This paper keeps the LazyFCA
candidate-generation mechanism but studies a different aggregation question:
after the local candidates have been generated, can they be ranked in a common pool and
reduced to a compact top-$k$ subset without losing predictive quality?

\section{Related Work}

This work is located at the intersection of FCA and pattern structures,
lazy premise-based classification, associative and rule-based classification,
and local interpretable machine learning. The common thread across these
areas is the tension between expressive local evidence and compact
human-readable decision support. FCA and pattern structures provide the
mathematical language for closed descriptions and interpretable premises;
associative classifiers and rule learners provide a long history of ranking,
pruning, and aggregating many candidate rules; local explanation methods
motivate query-dependent evidence; and prototype methods show that small sets
of source examples can be legitimate explanations rather than merely
compressed training data.

\subsection{FCA, Pattern Structures, and Concept-Based Learning}

Formal Concept Analysis was introduced as a lattice-theoretic framework for
organizing objects and attributes into formal concepts
\cite{wille2009restructuring}. The standard monograph by Ganter and Wille
\cite{ganter1999formal} gives the canonical definitions of formal contexts,
derivation operators, extents, intents, closure operators, and concept
lattices. In machine-learning terms, the appeal of FCA is that a concept is
not an opaque latent component: it is a maximal group of objects together
with the maximal description shared by that group. This makes concept-based
models attractive when the final decision must be accompanied by a symbolic
or set-theoretic explanation. Application-oriented treatments of concept data
analysis further emphasize the use of concept lattices for data analysis,
knowledge discovery, and navigation of structured descriptions
\cite{carpineto2004concept,ignatov2014introduction}.

Pattern structures generalize FCA from binary object-attribute tables to
complex descriptions equipped with a meet or similarity operation
\cite{ganter2001pattern}. This generalization is crucial for the present
paper because the target data are not purely binary. Numerical features are
represented by intervals, categorical features may be represented by sets or
one-hot Boolean attributes, and richer pattern structures can be defined for
graphs, sequences, relational descriptions, or logical formulas. Kuznetsov
\cite{kuznetsov2009pattern} presents pattern structures as a direct way to
process complex data without forcing all information through conceptual
scaling, noting that scaling can increase representation complexity and harm
the readability of the resulting structures. Later work discusses the role of
pattern structures in knowledge discovery and big data
\cite{kuznetsov2013fitting,kaytoue2015pattern}, numerical pattern mining
\cite{article}, relational and database-oriented pattern structures
\cite{baixeries2014characterization,kotters2020conjunctive}, and complex
sequential data \cite{buzmakov2016mining}. Projection-based work addresses
the computational burden of pattern structures by simplifying the description
space before or during mining \cite{buzmakov2015revisiting}; more recent
generalizations study broader algebraic settings such as pattern setups and
multistructures \cite{belfodil2020pattern}. These works establish the
mathematical and computational setting in which our interval-pattern
candidate classifiers live, but they do not directly answer which generated
local candidates should be retained for a specific prediction.

The concept-based machine-learning literature also studies how FCA
representations can support interpretable prediction. Kwuida and Ignatov
\cite{kwuida2020interpretability} discuss interpretability, similarity, and
attribute contribution in concept-based learning, including the role of
formal intents as classification hypotheses and the use of Shapley-style
attribution for hypotheses and stability indices. This is close in spirit to
our goal: after a concept-based or pattern-based classifier produces symbolic
evidence, one still needs to decide which parts of this evidence matter most
for a particular decision.

\subsection{Lazy Classification with Pattern Structures}

The most direct predecessor of the present paper is the work on lazy
classification with interval pattern structures for credit scoring
\cite{masyutin2015lazy}. That paper motivates pattern structures for
non-binary credit data, where full concept-lattice construction can be too
expensive, and distinguishes weak positive and negative premises or classifiers
from the stricter notion of weak hypotheses in an interval pattern-structure
setting. Its central observation is especially important for our study: in
high-dimensional numerical contexts, intersecting the test object with every
training object often produces very specific premises whose extent contains
only the test object and the source object. As a result, the number of
positive and negative premises can become close to the number of training
objects, and simple voting over such low-support premises becomes hard to
interpret and may be unreliable. To address this, Masyutin et al. propose a
stochastic modification that intersects the query with random subsamples of
class-specific contexts, controlled by subsample size, number of iterations,
and an alpha threshold.

Our work addresses the same failure mode from a different angle. Instead of
replacing object-query intersections by random batch-query intersections, we
keep the local query-source intersection mechanism and study whether the
generated candidates can be ranked and pruned after generation. This makes the
contribution complementary to randomized LazyFCA-like classification:
randomization tries to obtain broader premises by changing how candidates are
generated, whereas our ranking framework asks which already generated
candidates should be allowed to vote. The distinction matters experimentally
because a ranking method can be evaluated as a post-generation compactness
mechanism and can be compared against both all-candidate aggregation and
random top-$k$ selection.

The IPS-KNN method \cite{tomat2023interval} provides another closely related
but distinct direction. IPS-KNN uses interval pattern structures to make
distance-weighted $k$-nearest-neighbor classification interpretable. Rather
than producing a large set of LazyFCA candidate classifiers, it constructs an interval
pattern as a reason for classification from the query and a selected set of
nearest neighbors. This supports an important premise of the present paper:
a compact interval-pattern description, and in the limiting case even one
carefully chosen local classifier, can be a meaningful explanation. However,
IPS-KNN is not a direct baseline for our first experimental study because it
changes the underlying classifier to a nearest-neighbor method. We use it
primarily as conceptual evidence that compact local interval-pattern reasons
are useful and interpretable, while the present paper focuses on ranking and
pruning candidate classifiers within the LazyFCA family.

\subsection{Associative Classification, Rule Ordering, and Pruning}

Associative classification is a natural neighboring literature because it
also generates many interpretable class rules and then must decide which rules
to keep, order, and apply. The CBA framework of Liu, Hsu, and Ma
\cite{liu1998integrating} introduced classification based on class
association rules, where association-rule mining is restricted to rules whose
consequent is a class label. This line of work already contains two themes
that are central to our paper: mining may produce a large candidate rule set,
and predictive use requires selecting a manageable subset. CMAR
\cite{li2001cmar} extends this view by using multiple class-association rules
for prediction, pruning rules by confidence, correlation, and database
coverage, and showing that the single highest-confidence applicable rule is
not always the best basis for classification. CPAR \cite{yin2003cpar}
likewise emphasizes predictive rules and coverage analysis instead of merely
enumerating all frequent class-association rules.

Survey work on associative classification explicitly identifies rule
discovery, ranking, pruning, prediction, and evaluation as separate design
choices \cite{thabtah2007review}. This separation is useful for our paper
because LazyFCA candidate generation and LazyFCA candidate aggregation are
often discussed together, whereas our contribution isolates the aggregation
stage and studies the ranking function as the object of experimentation.
Wang, Xin, and Coenen \cite{wang2007novel} review common rule-ordering
mechanisms, including confidence-support orderings, weighted relative
accuracy, Laplace accuracy, and $\chi^2$-based orderings. Zaiane and Antonie
\cite{zaiane2005pruning} study pruning for associative classifiers and
motivate the problem in terms that map directly onto our setting: association
rule mining can produce so many human-readable rules that the model is no
longer human-readable, and pruning should reduce the rule set without
damaging accuracy. Their discussion of true positives and false positives of
individual rules is particularly close to our use of per-candidate
contingency values.

Lazy associative classification is also relevant. Baralis, Chiusano, and
Garza \cite{baralis2008lazy} argue that rich rule sets can improve
classification, but that pruning should remove harmful rules while preserving
useful evidence. Their L3 classifier uses lazy pruning and a compact rule-set
representation. Although this work is not based on FCA or pattern structures,
it is an important analogue: it treats the prediction-time use of a large
interpretable rule collection as a separate problem from rule generation.

\subsection{Rule-Quality and Candidate-Importance Measures}

The metrics evaluated in this paper draw from several traditions rather than
from a single established LazyFCA ranking theory. Precision and support are
standard in rule mining and associative classification
\cite{liu1998integrating,thabtah2007review}. Their limitations are also
well-known: very specific rules can have high precision but negligible
coverage, while broad rules can cover many objects without being especially
discriminative for a target class.

Weighted relative accuracy (WRAcc) is particularly relevant because it
explicitly combines relative accuracy with rule generality. Todorovski,
Flach, and Lavrac \cite{todorovski2000predictive} define WRAcc as
$p(B)(p(H \mid B)-p(H))$, where $B$ is the rule body and $H$ the target class.
They motivate it as an alternative rule-evaluation measure that trades off
improvement over the default class rule with the number of examples covered.
In our setting, this suggests that a useful LazyFCA candidate classifier should
not only be pure, but should also be nontrivial with respect to the class
prior and the number of objects it describes.

Subgroup discovery provides another source of quality measures. In that
literature, subgroups are interpretable descriptions evaluated by how their
target distribution differs from the reference population while also taking
subgroup size into account \cite{atzmueller2005exploiting}. This motivates
exceptionality-based metrics such as relative gain, binomial quality
functions, WRAcc-like scores, and other measures that penalize both trivial
coverage and weak deviation from the population. The connection to our work
is conceptual rather than algorithmic: each LazyFCA candidate can be viewed as
a local subgroup description induced by the query and a source object, and
ranking asks which of these descriptions are most informative for prediction.

FCA itself offers importance notions that are not purely predictive. Concept
stability measures how robust a concept intent is to removing objects from
its extent and has been used to reduce large concept lattices to more
meaningful taxonomies \cite{kuznetsov2009concept}. Stability is attractive
because it captures robustness rather than only accuracy, but it is expensive
and its meaning changes in our lazy setting: generated candidates are
query-dependent and often have very small extents. We therefore treat
stability-related quantities as one family of candidate metrics rather than
as the only principled importance criterion.

\subsection{Compact Interpretable Rule Models}

A broader interpretable machine-learning literature studies compact rule
sets, rule lists, and decision sets. Bayesian Rule Lists learn sparse ordered
if--then lists from pre-mined rules and use priors that favor short lists and
short rules \cite{letham2015interpretable}. Falling Rule Lists impose an
ordered risk structure on rule lists \cite{wang2015falling}. CORELS searches
for certifiably optimal sparse rule lists under a regularized empirical risk
\cite{angelino2018learning}. Bayesian Rule Sets similarly optimize sparse
sets of short rules for interpretable classification \cite{7837984}, and BETA
learns interpretable and explorable decision sets approximating black-box
models \cite{lakkaraju2017interpretable}. More recent work uses optimization
and satisfiability-based formulations to balance accuracy, rule size, and
scalability in rule-based classifiers \cite{ghosh2022efficient}.

These methods differ from our setting because they learn a global compact
model or a global surrogate. In contrast, LazyFCA produces query-specific
candidate classifiers, and the compact candidate set for one query need not be
the compact candidate set for another. Nevertheless, the rule-list and
rule-set literature supports a
key assumption of this paper: interpretability is not guaranteed merely
because individual rules are symbolic. The number of rules, their ordering,
their redundancy, and the way conflicts are resolved are part of the
interpretability problem.

\subsection{Locality, Similarity, and Prototype-Like Evidence}

Local explanation methods motivate the query-weighted metrics considered in
this paper. LIME explains an individual prediction by fitting a sparse
interpretable model locally around the query, using a proximity kernel to
weight perturbed samples \cite{ribeiro2016should}. The important idea for our
setting is not the perturbation mechanism itself, but the principle of local
fidelity: an explanation should be faithful in the neighborhood of the
instance being explained, and globally important features or rules may be
irrelevant locally. Query-similarity-weighted metrics adapt this principle to
pattern-structure candidates by favoring candidate descriptions that remain
close to the query while still providing class evidence.

Prototype and case-based methods give a second justification for local
evidence. Bien and Tibshirani \cite{bien2011prototype} study prototype
selection as interpretability through sparsity in samples rather than
sparsity in variables. ProtoPNet \cite{chen2019looks} shows the same idea in
a modern neural setting: predictions can be explained by similarity to learned
prototypes. These works are useful for interpreting singleton or near-singleton
LazyFCA candidates. In our setting, every generated candidate is constructed
from the query and at least one source object, so $tp \geq 1$ by construction.
A candidate with $tp=1$ should not be mistaken for a broadly supported formal
hypothesis or rule, but it can still carry local prototype-like evidence. The
empirical question is whether ranking functions can distinguish useful local
source-query evidence from harmful or redundant singleton candidates.

\subsection{Gap Addressed by This Work}

Existing work gives strong reasons to care about interpretable local
candidates, rule quality, pruning, compactness, and locality. However, the
specific problem studied here remains underexplored: given the large set of
query-dependent pattern-structure candidate classifiers produced by
LazyFCA-style classification, how should one rank candidates from all classes
in a single pool, and how many of the top-ranked candidates are actually
needed for an accurate and interpretable prediction? Randomized interval-pattern lazy
classification changes the generation process \cite{masyutin2015lazy};
associative classifiers prune and order global class-association rules
\cite{liu1998integrating,li2001cmar,baralis2008lazy}; compact rule learners
learn global rule models \cite{letham2015interpretable,angelino2018learning};
and local explanation methods explain black-box predictions
\cite{ribeiro2016should}. Our contribution is narrower and more specific:
we study post-generation ranking and compact top-$k$ selection for local
pattern-structure classifier candidates, with candidates from all classes
competing in one common ranking.

\section{Method}

\subsection{Global Pooled Top-k Selection}

Let $T=\{(x_i,y_i)\}_{i=1}^{n}$ be the training set, and let
$Y=\{0,\ldots,C-1\}$ be the set of class labels after encoding. For a query
object $q$ and a class $c \in Y$, LazyFCA generates one candidate classifier
from each training object of class $c$. We denote the resulting class-specific
candidate set by
\[
    \mathcal{P}_c(q) = \{z(q,s,c) \mid (s,y_s) \in T,\; y_s=c\}.
\]
The complete query-dependent candidate pool is
\[
    \mathcal{P}(q) = \bigcup_{c \in Y} \mathcal{P}_c(q).
\]
Each candidate $z \in \mathcal{P}(q)$ has a class label $y(z)$, an
interval-pattern description, contingency values $tp(z),fp(z),tn(z),fn(z)$, and optional
geometry or locality values derived from the description.

Given a ranking metric $m$, the implementation converts it into a score
$r_m(z)$ that is maximized. Most metrics are used directly. Metrics whose
natural interpretation is ``smaller is better'', such as false positives,
false negatives, error rate, and description volume, are multiplied by $-1$
for ranking. For an explanation budget $k$, the retained set is
\[
    \mathcal{P}_{k,m}(q) =
    \operatorname{top}_k\{z \in \mathcal{P}(q) \text{ sorted by } r_m(z)\}.
\]
If fewer than $k$ candidates are available, all available candidates are
retained.

The word ``global'' is important. The top-$k$ operation is applied to the
single pooled set $\mathcal{P}(q)$, not separately within each class. Therefore
the retained explanation may contain an unequal number of candidates per class.
This is intentional: a per-class quota can force weak candidates from a class to
remain, whereas a common pool lets the ranking metric suppress low-scoring
evidence regardless of class.

Prediction is based on a count vote over retained candidates. For each class
$c$, define
\[
    N_c(q) = |\{z \in \mathcal{P}_{k,m}(q) \mid y(z)=c\}|
\]
and
\[
    S_c(q) = \sum_{z \in \mathcal{P}_{k,m}(q):\,y(z)=c} r_m(z).
\]
The predicted class is selected lexicographically by retained count, summed
score, training-set class prior $\pi_c$, and finally the lower encoded label:
\[
    \hat{y}(q) =
    \operatorname*{arg\,max}_{c \in Y}^{\mathrm{lex}}
    \bigl(N_c(q), S_c(q), \pi_c, -c\bigr).
\]
The score sum is used only to resolve count ties; it does not replace the
primary voting rule. The research problem is to identify ranking metrics that
produce favorable accuracy-compactness tradeoffs: high macro-F1 with a small
number of retained interval-pattern candidates.

\section{Candidate Ranking Metrics}

The ranking metrics instantiate different notions of candidate importance:
purity, support, contrast against the class prior, statistical dependence,
geometric locality, and robustness. The implementation computes a broad set of
registered metrics, but the analysis will focus on the metrics that are both
interpretable and empirically useful.

For a candidate classifier $z$ predicting class $c$, let
\[
    p = tp(z)+fn(z), \qquad n = fp(z)+tn(z), \qquad M=p+n.
\]
Here $p$ is the number of supporters for class $c$, $n$ is the number of
opposers, and $M$ is the number of training objects. Let
\[
    \operatorname{cov}(z)=tp(z)+fp(z), \qquad
    \rho_c=\frac{p}{M}
\]
denote the total number of objects covered by the candidate and the training
prior of class $c$.

\subsection{Purity and Coverage Metrics}

The simplest metrics use the contingency table directly. Precision measures
the class purity of the covered objects:

\[
    \operatorname{precision}(z) = \frac{tp}{tp + fp}.
\]

Support measures the fraction of same-class training objects covered:

\[
    \operatorname{support}(z) = \frac{tp}{P},
\]

where $P=p$ is the number of training supporters for the candidate's class. The
implementation also records the false-positive rate over opposers,
$fp/n$, as an error-rate metric to be minimized, and the supporter-opposer
ratio $tp/fp$ with the usual infinite value when $fp=0$.

Lift compares the precision of the candidate with the class prior:
\[
    \operatorname{lift}(z) =
    \frac{\operatorname{precision}(z)}{\rho_c}.
\]
Values above one indicate that the covered region is more concentrated in
class $c$ than the training set as a whole. Support, confidence/precision, and
lift are standard rule-quality ingredients in associative classification and
rule mining \cite{liu1998integrating,thabtah2007review}.

The implementation also uses the following smoothed log-odds-like score:

\[
    \operatorname{LOR}(z) = \frac{2tp + 1}{2fp + 1}.
\]

This score is not a classical log odds ratio in logarithmic form. It is better
understood here as a support-sensitive purity score. When $fp=0$, it becomes
$2tp+1$, so it breaks ties among pure candidates by same-class coverage.
Two support-sensitive variants multiply precision or this odds-like score by
$\log(1+tp)$; another variant multiplies precision by $\sqrt{tp}$. These
variants test whether same-class coverage should be used as a soft tie-breaker
rather than as a hard filtering rule.

\subsection{Statistical and Rule-Learning Metrics}

Weighted relative accuracy (WRAcc) combines coverage with improvement over the
class prior:
\[
    \operatorname{WRAcc}(z) =
    \frac{tp+fp}{M}
    \left(\operatorname{precision}(z)-\rho_c\right).
\]
This follows the subgroup-discovery and predictive-rule-learning intuition
that a useful rule should cover a nontrivial region and should change the
target-class frequency relative to the default distribution
\cite{todorovski2000predictive,atzmueller2005exploiting}.

The implementation also includes several standard two-by-two table scores.
Youden's $J$ is
\[
    J(z)=\frac{tp}{tp+fn}-\frac{fp}{fp+tn},
\]
that is, sensitivity minus false-positive rate. Matthews correlation
coefficient is
\[
    \operatorname{MCC}(z)=
    \frac{tp\cdot tn - fp\cdot fn}
    {\sqrt{(tp+fp)(tp+fn)(tn+fp)(tn+fn)}}.
\]
These metrics treat each candidate as a one-versus-rest binary classifier
for its own predicted class.

Information gain and Gini gain evaluate how much splitting the training set
into covered and uncovered objects reduces label impurity. If
\[
    H(a,b)=-\frac{a}{a+b}\log_2\frac{a}{a+b}
           -\frac{b}{a+b}\log_2\frac{b}{a+b}
\]
is binary entropy, with zero terms omitted, then the implemented information
gain is
\[
    H(p,n) -
    \frac{(tp+fp)H(tp,fp)+(fn+tn)H(fn,tn)}{M}.
\]
Gini gain uses the same covered/uncovered decomposition but replaces entropy
by the binary Gini impurity
\[
    G(a,b)=1-\left(\frac{a}{a+b}\right)^2
             -\left(\frac{b}{a+b}\right)^2.
\]

Finally, the $\chi^2$ score and the likelihood-ratio $G$-test compare the
observed contingency table
\[
    \begin{pmatrix}
    tp & fp\\
    fn & tn
    \end{pmatrix}
\]
with the table expected under independence between coverage and class label.
For observed cells $O_i$ and expected cells $E_i$, the implemented scores are
\[
    \chi^2=\sum_i \frac{(O_i-E_i)^2}{E_i},
    \qquad
    G=2\sum_i O_i\log\frac{O_i}{E_i},
\]
with zero observed or expected terms handled safely.

\subsection{Query Locality and Pattern Geometry}

For interval pattern structures, geometry is part of interpretability. A
shorter interval pattern is usually easier to inspect and remains closer to
the query-source pair that generated it. For numerical interval patterns, the
implementation uses normalized interval widths

\[
    w_j(z) = \frac{\operatorname{upper}_j(z) - \operatorname{lower}_j(z)}
    {\operatorname{range}_j}.
\]

The interval tightness score is

\[
    \operatorname{tightness}(z) = 1 - \frac{1}{d}\sum_{j=1}^{d} w_j(z).
\]

Thus a degenerate interval around the query receives high tightness, while a
wide hyper-rectangle receives a lower score. Description volume is
\[
    \operatorname{volume}(z)=\prod_{j=1}^{d} w_j(z)
\]
and is minimized rather than maximized. The simplicity prior used by the
implementation penalizes description complexity through
\[
    \operatorname{simplicity}(z)=
    \frac{1}{1+\operatorname{complexity}(z)}.
\]

In the main numerical setting, query numeric similarity is the interval
tightness score, and the overall query similarity is the same value. When
one-hot Boolean features are present, query binary similarity is the fraction
of active query features retained in the candidate, and the overall query
similarity is the mean of the binary and numerical components. This gives a
simple locality score in the spirit of local explanation methods that weight
evidence by proximity to the explained instance \cite{ribeiro2016should}.

\subsection{Hybrid Locality-Purity Metrics}

The main hybrid metrics are:

\[
    \operatorname{QWP}(z) =
    \operatorname{precision}(z) \cdot \operatorname{query\_similarity}(z),
\]

and

\[
    \operatorname{QWLOR}(z) =
    \operatorname{LOR}(z) \cdot \operatorname{query\_similarity}(z).
\]

These metrics prefer candidates that are both class-pure and local to the
query. The working expectation is that such metrics better support compact
local explanations than global support alone.

The implementation also evaluates locality-weighted variants of the
support-sensitive scores, including precision multiplied by $\log(1+tp)$ or
$\sqrt{tp}$ and the locality-weighted WRAcc score. These hybrids are useful
because pure support can over-reward broad majority-class regions, while pure
locality can select intervals that are close to the query but weakly related
to the class label.

\subsection{Stability-Inspired Metrics}

FCA concept stability measures how robust a concept description is under
changes to the objects in its extent and is commonly used as an importance
notion for reducing large concept collections \cite{kuznetsov2009concept}.
The lazy setting is different from full lattice construction because each
candidate is query-dependent and may have a very small covered supporter
set. For this reason, the implementation treats stability-like quantities as
candidate ranking metrics rather than as definitive concept importance.

The implemented stability-inspired score is computed from witness counts for
the covered supporters: witnesses for retained or dropped Boolean conditions
and witnesses for the lower and upper numerical interval bounds. The
stability score multiplies factors of the form $1-2^{-w}$ over positive
witness counts $w$, and $\Delta$-stability is the minimum such witness count.
This gives a compact robustness proxy for local interval candidates, but its
interpretation remains tied to the lazy candidate-generation procedure.

\subsection{Singleton Source-Query Candidates}

A key methodological caveat is that every generated candidate is constructed
from the query and one same-class source object. Therefore, the source object is
covered by construction and $tp \geq 1$. A candidate with $tp=1$ is not invalid,
but it should be interpreted as a singleton source-query candidate rather than
as evidence of generalization to multiple same-class objects.

This distinction matters for both ranking and explanation. If a singleton
candidate also has $fp=0$, then precision is one, and many such candidates can
tie under pure contingency metrics. In that case, geometric locality,
interval tightness, or support-sensitive tie-breakers may determine which
source-query patterns are retained. Singleton candidates can still be useful:
they behave like local prototype evidence, connecting the query to a concrete
same-class training object through an interpretable interval pattern. The
paper therefore does not treat $tp=1$ as an automatic failure. Instead,
singleton rates are diagnostic quantities, and a possible $tp \geq 2$
restriction is best understood as an ablation or robustness check rather than
as the default definition of a valid LazyFCA candidate.

\section{Experimental Design}

\subsection{Datasets}

The main paper-level comparison is intended to focus on numerical datasets.
This keeps the interval pattern structure as the common description language
and makes the proposed ranked LazyFCA method comparable with the imported
FCALC, randomized FCALC, and IPS-KNN baselines, which are numerical-data
methods. The repository also contains a mixed churn dataset, but it is best
treated as a diagnostic or supplementary experiment unless the final paper
explicitly broadens beyond numerical data.

The numerical experimental suite contains:

\begin{itemize}[leftmargin=*]
    \item Breast cancer, binary, numerical after dropping the identifier column.
    \item Page blocks, multi-class, numerical.
    \item Parkinsons, binary, numerical.
    \item Rice, binary, numerical.
    \item Sonar, binary, numerical.
    \item Spambase, binary, numerical.
    \item Waveform, binary in the current local file, numerical.
    \item Ionosphere, binary, numerical.
    \item Image segmentation, multi-class, numerical.
    \item Vehicle, multi-class, numerical.
    \item Glass, multi-class, numerical after dropping the identifier column.
\end{itemize}

Final source citations, class distributions, and train/test sizes should be
reported after the experiment outputs are locked.

\subsection{Protocol}

For each dataset:

\begin{itemize}[leftmargin=*]
    \item Use stratified 80/20 train/test splits.
    \item Repeat over split seeds $1998,\ldots,2007$.
    \item Encode target labels as integers $0,\ldots,C-1$.
    \item For the ranked LazyFCA runs, pass numerical columns directly as
    interval features.
    \item Do not scale numerical features in the ranked LazyFCA runs, because
    interval widths and query similarities are computed relative to
    dataset-level numerical ranges.
    \item Use the same split seeds as the preserved imported baseline results
    so that macro-F1 summaries are comparable over aligned repeats.
\end{itemize}

\subsection{Compared Methods}

The paper compares the proposed selection mechanism with LazyFCA-family and
pruning baselines:

\begin{itemize}[leftmargin=*]
    \item Global top-$k$ ranked LazyFCA: all class-specific candidates are ranked in one
    pool and only top-$k$ are retained.
    \item Random top-$k$: candidates are randomly ordered as a sanity baseline.
    \item Deterministic FCALC/LazyFCA baseline: imported from the related
    interval-pattern lazy-classification study, with aggregation selected by
    cross-validation on the training split.
    \item Randomized FCALC/LazyFCA baseline: imported from the same study, with
    aggregation and sampling parameters selected by cross-validation on the
    training split.
\end{itemize}

The current raw-count all-candidate implementation is kept only as an
optional diagnostic. IPS-KNN is available as an optional external compact
interval-pattern baseline, but it changes the classifier family from
LazyFCA-style candidate aggregation to nearest-neighbor classification with a
single compact interval reason. If reported, it should therefore be framed as
family-level context rather than as the main baseline that defines success or
failure of the ranking method.

\subsection{Evaluation Measures}

The primary predictive metric is macro-F1 for both binary and multi-class
datasets. This choice is important because the imported FCALC, randomized
FCALC, and IPS-KNN result files contain macro-F1 over the aligned split seeds
but do not contain predictions or confusion matrices from which another
binary-class-specific score could be reconstructed.

The runner also stores accuracy, precision, recall, macro-F1, weighted-F1,
AUC-ROC when available, confusion matrices, mean retained candidates, and
compression ratio.

For compactness, the main summary reports:

\begin{itemize}[leftmargin=*]
    \item best macro-F1 and its corresponding $k$,
    \item smallest $k$ within 1 percent, 3 percent, and 5 percent of the best
    macro-F1,
    \item retained-candidate compression ratio.
\end{itemize}

\section{Implementation and Reproducibility}

The experiments are implemented in the repository under
\texttt{experiments/}. The runner is config-driven and incremental:

\begin{itemize}[leftmargin=*]
    \item \texttt{experiments/config.yaml} controls datasets, metrics,
    seeds, methods, and the $k$ grid.
    \item \texttt{experiments/run\_experiments.py} streams one
    test query at a time to avoid storing all explanations in memory.
    \item Result chunks are written per dataset, seed, method, and metric.
    \item Completed chunks are skipped on rerun.
\end{itemize}

This design was chosen because caching full LazyFCA explanations can require
excessive RAM and disk space on larger datasets.

\section{Results}

TODO: Fill after full experiment completion.

Suggested result tables:

\begin{table}[h]
\centering
\caption{Accuracy-compactness summary. TODO: Fill from compactness summary.}
\begin{tabular}{llrrrr}
\toprule
Dataset & Metric & Best $k$ & Best macro-F1 & Smallest $k$ within 3\% & Compression \\
\midrule
TODO & TODO & TODO & TODO & TODO & TODO \\
\bottomrule
\end{tabular}
\end{table}

\begin{table}[h]
\centering
\caption{Vanilla LazyFCA versus compact ranked LazyFCA. TODO: Fill from outputs.}
\begin{tabular}{lrrr}
\toprule
Dataset & Vanilla F1 & Best compact F1 & Compact $k$ \\
\midrule
TODO & TODO & TODO & TODO \\
\bottomrule
\end{tabular}
\end{table}

Suggested figures:

\begin{itemize}[leftmargin=*]
    \item One compactness-first top-$k$ curve per dataset.
    \item One plot per strong metric across datasets.
    \item A table or plot showing singleton $tp=1$ rates.
\end{itemize}

\section{Discussion}

TODO: Discuss which metric families work and why. Expected discussion points:

\begin{itemize}[leftmargin=*]
    \item Vanilla LazyFCA can be weak because equal aggregation over all
    candidates includes harmful candidates.
    \item Global ranking allows candidates from all classes to compete, so a
    query may retain an unbalanced but useful set of class evidence.
    \item Locality-aware purity metrics may be more robust for compact
    explanations than support multipliers.
    \item Metrics that directly reward support can be biased toward majority
    classes or broad rules in imbalanced datasets.
    \item Singleton candidates may carry useful nearest-neighbor-like local
    signal, but they should not be overinterpreted as generalized rules.
\end{itemize}

\section{Threats to Validity}

TODO: Discuss:

\begin{itemize}[leftmargin=*]
    \item Dataset suite size and representativeness.
    \item Sensitivity to train/test split seeds.
    \item Computational constraints and missing baselines.
    \item The difference between reducing explanation size and reducing
    candidate generation time.
    \item The interpretation of $tp=1$ source-query candidates.
\end{itemize}

\section{Conclusion}

TODO: Summarize the final findings. The intended conclusion is that ranked
global top-$k$ selection can make LazyFCA explanations dramatically more
compact, and in some cases improve predictive quality by avoiding harmful
candidate accumulation. The final statement must be grounded in the completed
experimental results.

\section*{Acknowledgments}

TODO: Add students, supervisors, grants, or lab acknowledgments if applicable.

\bibliographystyle{plain}
\bibliography{mybib}

\end{document}
